\section{Discussion}
\label{sec:discussion}

It is important to locate our negative result correctly: the limitation we establish is a property of the \emph{problem}, not of the decoupled fixed-point construction in particular. Because the conditional-mean ceiling of Proposition~\ref{prop:ceiling} applies to \emph{any} estimator that is a function of $\mathbf{x}_\mathrm{obs}$ and is trained under squared loss, it binds the direct regressor and any amortized point estimator exactly as it binds our iteration. Methods built on other principles fail as well, but in different ways. Likelihood-based estimators --- nonlinear least squares, or an EM scheme that alternates between inferring the latent trajectory and updating the parameter --- target a maximizer of the observation likelihood rather than a conditional mean; under non-identifiability the likelihood is flat along the fiber $\mathcal{F}_c$, so such methods converge to \emph{some} admissible fiber point, but a non-unique, initialization-dependent one, with no indication of the residual ambiguity. Thus squared-loss methods regress off the fiber to a biased mean, while likelihood-based methods land on the fiber but cannot single out the true solution; neither returns a correct, unique estimate, because the information required to do so is simply absent from $\mathbf{x}_\mathrm{obs}$.

This dichotomy points to what a satisfactory answer must look like. When a problem is non-identifiable, the object genuinely determined by the observation is not a single parameter but a set --- the fiber, together with the distribution induced on it by the prior and the observation noise. Any method that insists on returning one point must therefore either bias the estimate (squared loss) or make an arbitrary choice (likelihood), and no amount of algorithmic refinement removes this. The natural remedy is to lift inference from point estimates to \emph{distributions}: to characterize the posterior $p(\theta \mid \mathbf{x}_\mathrm{obs})$, which correctly represents the identifiable statistic ($mDI = S_I\sigma$ for OGTT), the residual uncertainty along the fiber, and the physical admissibility that point estimators forfeit. Developing such a distribution-valued amortized inference procedure --- one that retains the decoupled hidden-state structure but replaces the deterministic maps with conditional samplers --- is the direction we pursue in subsequent work, and the ceiling identified here is precisely what it is designed to overcome.
